\chapter{Introduction}
\label{ch:intro}

MolDStruct MC/MD simulates what happens to a biomolecule, or any small system of light atoms, when it is hit by an
intense femtosecond X-ray pulse.
The pulse ionizes the atoms, the released electrons partly escape and partly stay trapped, the sample charges up and
explodes, and the ions fly apart in directions that still carry information on the original structure.
The model follows all of this in one simulation: the electronic state of every atom, the electrons released, and the
nuclear motion.

It is part of the MolDStruct toolbox and is built on GROMACS~4.5.4~\cite{pronk2013}, extended with an ionization
module.
Everything GROMACS does still works the same way; the ionization model is switched on with \code{mdrun -ionize} and
controlled by \opt{mcmd-*} options in the \file{.mdp} file.

\section{The model in one page}

\begin{description}
  \item[Atoms] are classical particles that also carry an electronic configuration: the number of electrons in their
    K, L and M shells.
    Each configuration has a charge, and the charges enter the force field.
  \item[Electronic transitions] (photoionization by the pulse, Auger--Meitner decay, fluorescence) change those
    configurations.
    They are sampled every MD step by kinetic Monte Carlo from tabulated rates and cross sections
    (Chapter~\ref{ch:atomic}).
  \item[Charge transfer] moves electrons between nearby atoms when a classical over-the-barrier criterion allows it
    (Chapter~\ref{ch:ct}).
  \item[The force field] is modified so that ionized atoms lose their chemical bonds: bonded terms fade out with the
    charge, bonds break irreversibly, and the Coulomb interactions between bonded neighbours are restored
    (Chapter~\ref{ch:ff}).
  \item[Released electrons] either escape or are trapped by the growing positive charge.
    The trapped ones form a hot electron cloud, described as a continuous density in thermal equilibrium with the ions:
    a Poisson--Boltzmann problem solved on a three-dimensional grid every 0.2\,fs (Chapter~\ref{ch:pb}).
    The cloud screens the ions, pulls on them, decides which new electrons escape, and evaporates as the sample
    expands.
  \item[Electron--ion collisions] in the cloud ionize further (Chapter~\ref{ch:coll}).
  \item[Atomic data] (energies, rates and cross sections per configuration) are computed at start-up from a
    screened-hydrogenic shell model, or read from plain-text tables, which can be generated with XATOM (open source)
    or from CRETIN models (Chapter~\ref{ch:data}).
\end{description}

With \opt{mcmd-electrons = none} the electron cloud is switched off: every released electron leaves, and the model is
the original MolDStruct model~\cite{ostlin2018}.

\section{How this manual is organized}

Chapter~\ref{ch:start} gets a first run going.
Chapters~\ref{ch:atomic}--\ref{ch:coll} describe the physics, each with the options that control it.
Chapter~\ref{ch:data} specifies the atomic data, Chapter~\ref{ch:output} every output file.
Chapters~\ref{ch:mdp} and~\ref{ch:env} are the reference: every mdp option and every diagnostic environment variable.
The appendices cover the numerics of the electron model (Appendix~\ref{app:numerics}), parallel runs
(Appendix~\ref{app:parallel}), validation and limitations (Appendix~\ref{app:validation}), and units, the source
files and the version history (Appendix~\ref{app:misc}).

Basic knowledge of GROMACS is assumed; see the GROMACS documentation for anything not specific to the ionization
model.

\section{Conventions}

Inside the model energies are in eV, lengths in nm, and times in ps as in GROMACS --- except the pulse options and
the electron-model intervals, which are in \textbf{femtoseconds}.
Charges are in units of $e$.
A configuration is written \code{K L M}, e.g.\ \code{2 4 0} for neutral carbon and \code{1 4 0} for carbon with a
K-shell hole.
Potentials are \emph{potential energies of an electron}, in eV.

\section{Citing}

If you use MolDStruct MC/MD, please cite the MolDStruct paper (to be added on publication), and GROMACS~4.5
\cite{pronk2013}.
The built-in atomic data follow Scott and Hansen~\cite{scott2010}; if you use atomic data generated with XATOM, cite
XATOM~\cite{jurek2016}.
